\documentclass{article}
\usepackage[margin=1.25in]{geometry}

\usepackage{graphicx}
\usepackage{amsmath,amssymb,amsthm}
\usepackage{stmaryrd}
\usepackage{xcolor}
\usepackage{enumerate}

\usepackage[T1]{fontenc}
\usepackage{stix2}

% \renewcommand{\familydefault}{\sfdefault}

\usepackage[hidelinks]{hyperref}

\theoremstyle{plain}
\newtheorem{lemma}{Lemma}
\newtheorem{theorem}{Theorem}
\newtheorem{proposition}{Proposition}
\newtheorem{assumption}{Assumption}
\newtheorem{definition}{Definition}
\newtheorem{corollary}{Corollary}
\newtheorem{example}{Example}

\theoremstyle{remark}
\newtheorem{remark}{Remark}

\setlength{\parskip}{0.5em}
\setlength{\parindent}{0em}

\newcommand{\Dfx}{Df}

\input{macros}

\title{Notes on the Mixed Jacobian Tensor}
\author{Akash Harapanhalli}

\begin{document}
\maketitle

\section{Introduction}
In this note I present a higher-order analogue of the mixed Jacobian approach.

To do so, we will define a new \emph{mixed }
mixed evaluation approach towards bounding a functions behavior

\section{Some Background}

\subsection{The Mixed Jacobian Matrix/Operator}

Recall the definition of the mixed Jacobian operator
\begin{align*}
    (M_{x'}f(x,s))_{ij} = \frac{\partial f_i}{\partial x_j} (x_1,\dots,x_{j-1},s_jx_j + (1-s_j)x'_j,x'_{j+1},\dots,x'_n),
\end{align*}
which coupled with $n$ mean value evaluations along the element-wise path connecting $x'$ to $x$ results in the following,
\begin{align*}
    f(x) = f(x') + M_{x'}f (x,s) (x - x'),
\end{align*}
for some $s\in[0,1]^n$. 

\subsection{Some Tensor Notation for Higher Order Derivatives}

Let $f:\R^n\to\R^m$ be a smooth function. We know that $Df(x)$ denotes the Jacobian matrix, which is an operator mapping vectors $v \mapsto Df(x)[v]$ to their directional derivatives through $f$. 
We can of course iterate this map to obtain, e.g., $D^kf$, which maps a collection of vectors $[v_1,\dots,v_k] \mapsto $ to the higher order directional derivative of $f$ in the corresponding direction (``how the change in $v_1$ is changing as a function of directional change in $v_2$, etc...'').

The best way to think of $D^kf(x)$ is as a symmetric tensor mapping $k$ elements of $T_x\R^n$ to $T_{f(x)}\R^m$. 
Computationally, we can think of $D^kf(x)$ as a multi-dimensional array defined by the following: for any sequence of indices $\bfj = (j_1,\dots,j_k) \in \{1,\dots,n\}^k$, 
let $\alpha(\bfj)$ be the corresponding multi-index defined by counting the number of occurrences of each index,
\begin{align*}
    \alpha(\bfj)_i = |\{\ell\in\{1,\dots,k\} : j_\ell = i\}|
\end{align*}
Then the multi-dimensional array is defined by
\begin{align*}
    (D^kf(x))_{\bfj} 
    = \frac{\partial^k f}{\partial x_{j_1} \cdots \partial x_{j_k}} 
    = \frac{\partial^k f}{\partial x^{\alpha(\bfj)}}(x) = \frac{\partial^k f}{\partial x^{\alpha(\bfj)_1}_1 \cdots\partial x^{\alpha(\bfj)_n}_n} (x),
\end{align*}
which is often a more useful definition for working computationally.
The multi-index simplification is due to Schwarz's theorem---since $D^kf$ is symmetric, the multi-dimensional array is redundant, in the sense that any two index sequences $\bfj_1, \bfj_2$ where $\alpha(\bfj_1) = \alpha(\bfj_2)$ will satisfy $(D^kf(x))_{\bfj_1} = (D^kf(x))_{\bfj_2}$.

For instance, JAX will build this automatically through chained \verb|jacfwd| calls: $k$ iterations of \verb|jacfwd| evaluated at \verb|x| will return an array of shape \verb|(m,n,...,n)| (with $k+1$ dimensions). 

We often use $D^kf$ to perform a Taylor series expansion, which using this notation can be written as
\begin{align*}
    f(x) = f(x') + \sum_{k=1}^p \frac1{k!} D^kf(x') [(x - x')^{\otimes k}] + O(|x - x'|^{p+1}),
\end{align*}
using tensor contraction. 
Using multi-index notation this sum can be expressed simpler as
\begin{align*}
    f(x) = f(x') + \sum_{|\alpha| \leq p} \frac{1}{\alpha!} \frac{\partial^{|\alpha|}f}{\partial x^\alpha} (x - x')^\alpha + O(|x - x'|^{p+1}).
\end{align*}
This illustrates how tensor contraction with the dense $D^kf$ is often overkill---it individually sums over each of the index sequences $\bfj$ that correspond to the same multi-index $\alpha(\bfj)$.
Combinatorially, given a multi-index $\alpha$ we can compute the number of corresonding index sequences $\bfj$: this is given by the multinomial coefficient $|\alpha|!/(\alpha_1!\alpha_2!\cdots\alpha_n!)$, compactly written as $|\alpha|!/\alpha!$.
This explains why the inverse factorial in the second expression is $1/\alpha!$ instead of $1/|\alpha|!$.

Rather than working with the full dense tensor $D^kf$, we can instead work with a \emph{lower triangular} tensor $L^kf$ called the \emph{reduced derivative tensor},
\begin{align*}
    (L^kf (x))_{\bfj} = \begin{cases}
        \displaystyle \frac{1}{\alpha(\bfj)!} \frac{\partial^{k} f}{\partial x^{\alpha(\bfj)}}(x) & j_1 \geq j_2 \geq \cdots \geq j_k \\
        0 & \text{otherwise}
    \end{cases}.
\end{align*}
Thus we can rewrite, once again, the Taylor expansion to the following,
\begin{align*}
    f(x) = f(x') + \sum_{k=1}^p L^kf(x')[(x - x')^{\otimes k}] + O(|x - x'|^{p+1}).
\end{align*}
While this tensor contraction seems just as complex as before, there is enough sparsity to exploit in the lower triangular structure $L$ to make it equivalent to the second expression with multi-indices.

One way to think about the reduced tensor $L^kf$ is that it is essentially indexed by multi-indices themselves, since each nonzero entry in $L^kf$ is in one-to-one correspondence with a multi-index $\alpha$. 
We can see this by mapping $\alpha$ to the unique sorted element $\bfj$ in the preimage of $\alpha(\cdot)^{-1}(\alpha)$ to obtain
\begin{align*}
    (L^kf)_{\alpha} = (L^kf)_{\bfj}.
\end{align*}

% We can turn a multiindex $\alpha$ into one choice of index by choosing one variable at a time and decrementing the corresponding index. For example, $(2,3,1)$ can be represented by \verb|[0,0,1,1,1,2]| or \verb|[1,0,1,0,2,1]|. 

% Instead of iterated \verb|jacfwd|, we might consider using \verb|jvp|, or even better, \verb|jet| to extract all of these through one call.
% To do so, define the intermediate function
% \begin{align*}
%     g(s) = f(x' + s(x - x')),
% \end{align*}
% noting that $\frac{\rmd}{\rmd t^k} g(s) = D^kf(x')[(x - x')^{\otimes k}]$.
% The function \verb|jet|, with the primal input \verb|x'| and the series input \verb|[(x - x'), 0, ..., 0]| will automatically extract all of these coefficients.

\newpage

\section{The Mixed Derivative Interval Tensor}

\begin{definition}[The Mixed Derivative Interval Tensor]
Let $f:\R^n\to\R^m$ be $C^k$ smooth. 
Let $X\in\IR^n$, and $x'\in X$ be a fixed point. 
The \emph{mixed derivative interval tensor} $M^k_{x'}f(X) : \underbrace{\R^n \times \cdots \times \R^n}_{k\text{ copies}} \to \IR^m$ is defined in the following way.
Define a lower triangular multi-dimensional interval array as follows,
\begin{align*}
    (M^k_{x'}f (X))_{\bfi} = \begin{cases}
        \displaystyle \frac{1}{k\alpha(\bfi)!} \sum_{j=1}^n \alpha(\bfi)_j \frac{\partial^{k} f}{\partial x^{\alpha(\bfi)}} (X_1,\dots,X_j,x'_{j+1},\dots,x'_n) & i_1 \geq i_2 \geq \cdots \geq i_k \\
        0 & \text{otherwise}
    \end{cases}.
\end{align*}
The contraction of $M^k_{x'}f(X)$ with a tensor power $v^{\otimes k}$ results in an interval in $\IR^m$ defined as follows,
\begin{align*}
    M^k_{x'}f (X) [v^{\otimes k}] = \sum_{|\alpha| = k} (M^k_{x'}f(x))_{\bfi(\alpha)} v^\alpha
\end{align*}
where a multi-index $\alpha$ once again corresponds uniquely with the unique sorted element $\bfi(\alpha)$.
\end{definition}

There does exist a mean value form of the Theorem, however it is quite complex and I'm not sure we need it. 
It involves a different $s_\alpha\in(0,1)^n$ for each multi-index $\alpha$, so many mean value choices. But there is an equality with that route.

\begin{lemma} \label{lem:interval_integ}
Let $[\ula,\ola]$ be a scalar interval, and $p:[t_0,t_f]\to\R_{\geq0}$ be a nonnegative function. Then the following holds, 
\begin{align*}
    \left\{\int_{t_0}^{t_f} a(t) p(t) \ \rmd t : a:[t_0,t_f]\to[\ula,\ola] \text{ measurable} \right\} 
    = [\ula,\ola] \int_{t_0}^{t_f} p(t) \ \rmd t.
\end{align*}
\end{lemma}
\begin{proof}
Since $p(t)\geq 0$ and $\ula \leq a(t) \leq \ola$, we have $\ula p(t) \leq a(t)p(t) \leq \ola p(t)$.
We have $\subseteq$ due to monotonicity of the Lebesgue integral.
Finally, the choice $a(t) = \ula$ and $a(t) = \ola$ prove $\supseteq$.
\end{proof}

\begin{theorem}
Let $f:\R^n\to\R^m$ be a $C^{p+1}$ smooth function for some $p\geq 0$, $X\in\IR^n$ be an interval, and consider some fixed $x'\in X$.
The following inclusion holds:
\begin{align*}
    f(x) &\in f(x') + \sum_{i=1}^p L^i f(x') [(x - x')^{\otimes i}] + M^{p+1}_{x'} f (X) [(x - x')^{\otimes (p+1)}] \\
    &= f(x') + \sum_{1\leq |\alpha|\leq p} \frac{1}{\alpha!} \frac{\partial^{|\alpha|} f}{\partial x^\alpha} (x') (x - x')^{\alpha} 
    + \sum_{|\beta| = p+1} \sum_{j=1}^n \frac{\beta_j}{(p+1)\beta!} \frac{\partial^{|\beta|} f}{\partial x^\beta} (X_1,\dots,X_j,x'_{j+1},\dots,x'_n) (x - x')^\beta
\end{align*}
\end{theorem}
\begin{proof}
($m=1$) We start by proving the scalar function case.
Applying Taylor's theorem in integral form along the line segment $tx + (1-t)x'$, 
\begin{align*}
    f(x) = f(x') + \sum_{|\alpha|\leq p-1} \frac{1}{\alpha!} \frac{\partial^{|\alpha|} f}{\partial x^\alpha} (x') (x - x')^{\alpha} + \sum_{|\alpha|=p} \frac{p}{\alpha!} (x-x')^\alpha \int_0^1 (1-t)^{p-1} \frac{\partial^pf}{\partial x^\alpha} (tx + (1-t)x')\ \rmd t.
\end{align*}
Fix a multi-index $|\alpha| = p$. Applying the mixed Jacobian for fixed $t$, there is a $s_{\alpha,t} \in (0,1)^n$ such that
\begin{align*}
    \frac{\partial^pf}{\partial x^\alpha} (tx + (1-t)x') 
    = \frac{\partial^pf}{\partial x^\alpha} (x') + t \sum_{j=1}^n \frac{\partial}{\partial x_j} \frac{\partial^pf}{\partial x^\alpha} (\gamma_j(tx + (1-t)x', x', s_{\alpha,t,j})) (x_j - x'_j),
\end{align*}
since $tx + (1-t)x' - x' = t(x - x')$. 
Plugging this back into the integral, the first term results in
\begin{align*}
    \int_0^1 (1 - t)^{p-1} \frac{\partial^pf}{\partial x^\alpha} (x') \ \rmd t = 
    \frac{\partial^pf}{\partial x^\alpha} (x') \int_0^1 (1 - t)^{p-1} \ \rmd t = \frac1p \frac{\partial^p f}{\partial x^\alpha} (x')
\end{align*}
which combines with the first sum, simply changing its bound to $|\alpha| \leq p$. 

Swapping the sum with the integral, the second term is
\begin{align*}
    I = \int_0^1 p (1 - t)^{p-1} t \sum_{|\alpha| = p} \sum_{j=1}^n \frac{1}{\alpha!}  \frac{\partial}{\partial x_j} \frac{\partial^pf}{\partial x^\alpha} (\gamma_j(tx + (1-t)x', x', s_{\alpha,t,j})) (x_j - x'_j)(x - x')^{\alpha}  \ \rmd t
\end{align*}
For each pair $(\alpha,j)$, associate the order $(p+1)$ multi-index $\beta$ by $\beta = \alpha + e_j$, noting of course that $\frac{\partial}{\partial x_j} \frac{\partial^p}{\partial x^\alpha} = \frac{\partial^{p+1}}{\partial x^\beta}$ and $(x_j - x'_j)(x - x')^{\alpha} = (x - x')^\beta$. Note that
\begin{align*}
    \frac{1}{\alpha!} = \frac{1}{\beta_1!\cdots\beta_{j-1}!(\beta_j-1)!\beta_{j+1}!\cdots\beta_n!} = \frac{\beta_j}{\beta!}.
\end{align*}
Next, for a particular $(\alpha,j)$ mapping to $\beta$, we have
\begin{align*}
    \frac{1}{\alpha!}  \frac{\partial}{\partial x_j} \frac{\partial^pf}{\partial x^\alpha} (\gamma_j(tx + (1-t)x', x', s_{\alpha,t,j}))  
    &= \frac{\beta_j}{\beta!} \frac{\partial^{p+1}f}{\partial x^\beta} (\gamma_j(tx + (1-t)x', x', s_{\alpha,t,j})) 
    \\
    &\in \frac{\beta_j}{\beta!} \frac{\partial^{p+1} f}{\partial x^\beta} (X_1,\dots,X_j,x'_{j+1},\dots,x'_n).
\end{align*}
Thus, we can replace the sum over $p$ multi-indices to sum over $(p+1)$ multi-indices as
\begin{align*}
    I \in \int_0^1 p (1 - t)^{p-1} t \underbrace{\sum_{|\beta| = p+1} \sum_{j=1}^n \frac{\beta_j}{\beta!} \frac{\partial^{p+1} f}{\partial x^\beta} (X_1,\dots,X_j,x'_{j+1},\dots,x'_n)(x - x')^\beta}_{M^{p+1}_{x'}f(X)[(x - x')^{\otimes (p+1)}]}   \ \rmd t.
\end{align*}
We interpret this inclusion in the nonconstant sense, meaning there is a real choice of measurable function $A(t)\in M^{p+1}_{x'}f(X)[(x - x')_\beta]$ such that $I = \int_0^1 p (1 - t)^{p-1} t A(t) \ \rmd t$.

Since the evaluation $M^{p+1}_{x'}f(X)[(x - x')^{\otimes (p+1)}]$ is simply an interval, and $(1 - t)^{p-1} t$ is nonnegative over the integration domain $[0,1]$, 
we can apply Lemma~\ref{lem:interval_integ} to pull the interval outside the integral.
Finally, since $\int_0^1 p(1-t)^{p-1} t \ \rmd t = \frac{1}{p+1}$, we obtain
\begin{align*}
    I\in M_{x'}^{p+1}f(X)[(x - x')^{\otimes (p+1)}] = \frac{1}{p+1}  \sum_{|\beta| = p+1} \sum_{j=1}^n \frac{\beta_j}{\beta!} \frac{\partial^{p+1} f}{\partial x^\beta} (X_1,\dots,X_j,x'_{j+1},\dots,x'_n) (x - x')^\beta.
\end{align*}

($m>1$) Let $\ell\in\R^m$ and apply the scalar case to the function $\ell^T f$. Optimizing over the set inclusion on the RHS, we see that
\begin{align*}
    \ell^T f(x) - \ell^T \left(f(x') + \sum_{i=1}^p L^i f(x') [(x - x')^{\otimes i}] \right) \leq \sup_{v\in M^{p+1}_{x'} f (X)[(x - x')^{\otimes (p+1)}]} \ell^T v.
\end{align*}
Thus, $f(x) - \left(f(x') + \sum_{i=1}^p L^i f(x') [(x - x')^{\otimes i}] \right)$ belongs to every halfspace containing the set $M^{p+1}_{x'} f (X)[(x - x')^{\otimes (p+1)}]$. 
As an interval subset, this set is closed and convex and thus it is characterized by these halfspaces, meaning
\begin{align*}
    f(x) - \left(f(x') + \sum_{i=1}^p L^i f(x') [(x - x')^{\otimes i}] \right) \in M^{p+1}_{x'} f (X) [(x - x')^{\otimes (p+1)}],
\end{align*}
completing the proof.
\end{proof}

% Every $\beta$ with $\beta_j\neq 0$ has an order $p$ multi-index 
% \begin{enumerate}[i)]
%     % \item (Mean value equality form) There exists $s_0\in(0,1)$ and $s_{\alpha}\in(0,1)^n$ for each multi-index $|\alpha|=p$ such that
%     % \begin{align*}
%     %     % f(x) = f(x') + \sum_{i=1}^{p} L^i f(x') [(x - x')^{\otimes i}] + s_0 \sum_{|\beta|=p+1} \frac{1}{\beta!} \sum_{j=1}^n \beta_j \frac{\partial^{p+1}f}{\partial x^\beta}(\gamma_j(s_0x + (1-s_0)x'),x',s_j) (x - x')^\beta
%     %     f(x) &= f(x') + \sum_{|\alpha| \leq p} \frac{1}{\alpha!} \frac{\partial^{|\alpha|} f}{\partial x^\alpha} (x') (x - x')^{\alpha} \\
%     %     & \ + s_0 \sum_{|\alpha|=p+1} \frac{1}{\alpha!} \sum_{j=1}^n \alpha_j \frac{\partial^{p+1}f}{\partial x^\alpha}(\gamma_j(s_0x + (1-s_0)x'),x',s_j) (x - x')^\alpha
%     % \end{align*}
%     \item (Tensor set inclusion form) The following inclusion holds,
% \end{enumerate}
% (Part i) Applying Taylor's theorem in mean value form along the line segment $sx + (1-s)x'$, we first see that there must exist a $s_0\in(0,1)$ such that
% \begin{align*}
%     f(x) = f(x') + \sum_{|\alpha|\leq p-1} \frac{1}{\alpha!} \frac{\partial^{|\alpha|} f}{\partial x^\alpha} (x') (x - x')^{\alpha} + \sum_{|\alpha|=p} \frac{1}{\alpha!} \frac{\partial^pf}{\partial x^\alpha} (s_0x + (1-s_0)x') (x - x')^\alpha.
% \end{align*}
% Next, fix a multi-index $|\alpha| = p$. Applying the mixed Jacobian operator to $\frac{\partial^p f}{\partial x^\alpha}$, there exists a $s_\alpha\in(0,1)^n$ such that
% \begin{align*}
%     \frac{\partial^p f}{\partial x^\alpha} (s_0x + (1-s_0)x') = \frac{\partial^p f}{\partial x^\alpha} (x') + \sum_{j=1}^n \frac{\partial}{\partial x_j} \frac{\partial^p f}{\partial x^\alpha} (\gamma_j(s_0x + (1 - s_0)x'),x',s_{\alpha,j}) (s_0(x - x')).
% \end{align*}

\begin{remark}[Intuitions]
What does the sum of $\alpha(\bfi)_j$ mean, and the division by $k$?
Suppose we consider the non-mixed pure interval version of this tensor replacing the mixed evaluations with interval evaluations. We end up with the following,
\begin{align*}
    (M^k_{x'}f (X))_{\bfi} = \begin{cases}
        \displaystyle \frac{1}{k\alpha(\bfi)!} \sum_{j=1}^n \alpha(\bfi)_j \frac{\partial^{k} f}{\partial x^{\alpha(\bfi)}} (X_1,\dots,X_j,X_{j+1},\dots,X_n) & i_1 \geq i_2 \geq \cdots \geq i_k \\
        0 & \text{otherwise}
    \end{cases}.
\end{align*}
However, the evaluation $\frac{\partial^{k} f}{\partial x^{\alpha(\bfi)}} (X)$ is not even dependent on the summation index $j$ anymore! Pull it out, and we are left with a $\sum_{j=1}^n \alpha(\bfi)_j$ which is just equal to $k$ since $\alpha$ is a $k$ multi-index. This $k$ clears the $k$ in the denominator and we are left with the usual $\frac{1}{\alpha(\bfi)!}$.

So what is the $\alpha(\bfi)_j$ doing in the mixed case? It is performing a partition of unity-like \emph{weighted} sum over the partial derivative contributions. 
Something like the following: along the $j$-th path, we peeled off $\alpha(\bfi)$ partial derivative layers with the same mean value application, and therefore need to scale the contribution accordingly.

There is another perspective we can analyze for a moment. Suppose we swap the sum over the multi-indices and the sum over $j$ to obtain something like the following:
\begin{align*}
    \frac{1}{(p+1)}\sum_{j=1}^n \sum_{|\beta| = p+1} \frac{\beta_j}{\beta!} \frac{\partial^{|\beta|} f}{\partial x^\beta} (X_1,\dots,X_j,x'_{j+1},\dots,x'_n) (x - x')^\beta
\end{align*}
The $\frac{\partial^{|\beta|}f}{\partial x^\beta}$
This interpretation is somehow building the path integral piece by piece. Step along each path, and catalog the partial derivative contributions \emph{along that piece of path}, meaning we only need to input $X_1,\dots,X_j,x'_{j+1},\dots,x'_n$. 
\end{remark}


\newpage

\bibliographystyle{ieeetr}
\bibliography{references.bib}

\end{document}

